\section{Introduction}

Digital patient frameworks promise computable representations of human biology that can support prediction, personalization, simulation, education, and clinical decision-making. The terminology is broad: healthcare reviews use labels including digital twin, patient digital twin, virtual patient, in silico patient, synthetic patient, digital avatar, and digital patient, often with overlapping but non-identical meanings.

This conceptual spread matters because different terms imply different technical commitments. A virtual or in silico patient may be static, archetype-based, or population-level, whereas the term digital twin often implies tighter patient-specific and temporal coupling. Without clearer boundaries, claims about maturity, validation, and clinical readiness remain difficult to compare.

Existing standards and guidance address complementary parts of this problem rather than a single end-to-end framework for digital patients. Cross-domain consensus characterizes a digital twin by dynamic updating, predictive capability, decision relevance, and bidirectional interaction with its physical counterpart, while regulatory credibility frameworks make the evidentiary burden dependent on a defined context of use and model risk.\cite{nasem2024digitaltwins,fda2023credibility,asme2018vv40} Construct identity should therefore be distinguished from fitness for a clinical purpose.

Guided by the PROSPERO protocol, this meta-review addresses six prespecified review questions: conceptual framing; construction and modeling frameworks; validation and fidelity; the application spectrum and functional roles; barriers and infrastructural gaps; and future directions and standardization priorities. The data-extraction strategy is viewpoint-driven rather than purely bibliographic. Each coded field populates a specific interpretive layer: study characteristics define the evidence base; terminology clarifies conceptual boundaries; hierarchy, data coupling, and temporality locate constructs on the digital-patient spectrum; modeling inputs and outputs describe the computational pipeline; validation fields address credibility and clinical readiness; application fields define the use-case landscape; and barrier and future-direction fields inform an implementation roadmap.

The objective of this manuscript is to clarify how digital patient constructs are defined, built, validated, and applied across recent healthcare review literature, while making the descriptive synthesis fully reproducible from the accompanying extraction workbook.
